% !TEX program = pdflatex
% !TEX encoding = UTF-8 Unicode
\documentclass[12pt,a4paper]{article}
\usepackage[utf8]{vietnam}
\usepackage{amsmath,amssymb}
\usepackage{tabularx}
\usepackage{tikz}
\usepackage{xcolor}
\usepackage[top=0.45cm, bottom=0.45cm, left=0.6cm, right=0.6cm]{geometry}
\usepackage[dethi]{ex_test}

\setlength{\parskip}{0pt}
\setlength{\parindent}{0pt}

% Dòng kẻ chấm mờ đều đặn chuẩn viết tay (khoảng cách 19.5pt)
\newcommand{\dongke}[1]{%
  \begingroup%
  \setlength{\parskip}{0pt}%
  \baselineskip=19.5pt%
  \vspace{1mm}%
  \foreach \i in {1,...,#1}{%
    \noindent{\color{gray!60}\dotfill}\par%
  }%
  \vspace{0.8mm}%
  \endgroup%
}

% Định nghĩa ô bài tập tự luận trong bảng
\newcommand{\oBaiTap}[3][5]{%
  \vspace{0.8mm}%
  \noindent\textbf{#2} #3\par%
  \vspace{0.8mm}%
  \noindent\textit{Bài giải:}%
  \dongke{#1}%
}

% Câu trắc nghiệm 4 đáp án dàn hàng ngang
\newcommand{\cauTN}[6]{%
  \noindent\textbf{Câu #1.} #2\par\vspace{1mm}%
  \noindent%
  \makebox[0.25\linewidth][l]{\textbf{A.} #3}%
  \makebox[0.25\linewidth][l]{\textbf{B.} #4}%
  \makebox[0.25\linewidth][l]{\textbf{C.} #5}%
  \makebox[0.25\linewidth][l]{\textbf{D.} #6}\par\vspace{1.5mm}%
}

\begin{document}
\pagestyle{empty}


% #############################################################################
% ========================= TOÁN 8 - TUẦN 03 - BUỔI 01: MÃ ĐỀ 101 ============
% ========================= MẶT TRƯỚC (TRANG 1) ===============================
% #############################################################################

\noindent
\begin{tabularx}{\linewidth}{|X|c|}
\hline
\textbf{PHIẾU KIỂM TRA ĐỊNH KỲ TOÁN 8 - TUẦN 03 - BUỔI 01} & \textbf{Mã đề: 101} \\
\textit{Chủ đề: Phép chia đa thức cho đơn thức} & \textit{Thời gian: 45 phút} \\
\hline
\end{tabularx}

\vspace{1mm}

\noindent
\begin{tabularx}{\linewidth}{|X|}
\hline
\rule[-2mm]{0pt}{7.5mm}\textbf{Họ và tên:} \textbf{\large Chu Gia Huy} \hfill \textbf{Mã HS:} \texttt{GIAHUY0503} \quad \\
\hline
\end{tabularx}

\vspace{1mm}

\noindent
\begin{tabularx}{\linewidth}{|p{3.2cm}|X|}
\hline
\centering\textbf{ĐIỂM SỐ} & \centering\textbf{LỜI NHẬN XÉT CỦA THẦY CÔ} \tabularnewline
\hline
\rule{0pt}{1.3cm} & \tabularnewline
\hline
\end{tabularx}

\vspace{1.2mm}

\noindent\textbf{BẢNG ĐIỀN ĐÁP ÁN TRẮC NGHIỆM (Mỗi câu đúng 0,5 điểm):}
\vspace{0.8mm}

\noindent
\renewcommand{\arraystretch}{1.35}
\begin{tabularx}{\linewidth}{|c|*{5}{>{\centering\arraybackslash}X|}}
\hline
\textbf{Câu} & \textbf{1} & \textbf{2} & \textbf{3} & \textbf{4} & \textbf{5} \\
\hline
\textbf{Đáp án} & & & & & \\
\hline
\end{tabularx}
\renewcommand{\arraystretch}{1.0}

\vspace{1.5mm}

\noindent\textbf{\large PHẦN I. TRẮC NGHIỆM KHÁCH QUAN (2,5 điểm)}
\vspace{1mm}

\cauTN{1}{Kết quả của phép chia đơn thức $15x^5 y^3 : 3x^2 y$ là:}{$5x^3 y^2$}{$5x^7 y^4$}{$12x^3 y^2$}{$5x^3 y$}

\cauTN{2}{Đơn thức $12x^4 y^3$ chia hết cho đơn thức nào dưới đây?}{$4x^5 y^2$}{$3x^3 y^4$}{$-6x^2 y^3$}{$2x^4 y^5$}

\cauTN{3}{Kết quả phép chia $(12x^4 y^2 - 8x^3 y + 4x^2) : 4x^2$ là:}{$3x^2 y^2 - 2xy$}{$3x^2 y^2 - 2xy + 1$}{$3x^2 y^2 - 2x + 1$}{$3x^6 y^2 - 2x^5 y + x^4$}

\cauTN{4}{Đa thức $P = 6x^3 y^2 - 9x^2 y^3 + 15xy$ chia hết cho đơn thức nào sau đây?}{$3xy^2$}{$3x^2 y$}{$5x^2 y^2$}{$-3xy$}

\cauTN{5}{Biết $(24x^4 y^3 - 16x^3 y^2) : A = 6x^2 y^2 - 4xy$. Đơn thức $A$ là:}{$4xy$}{$4x^2 y$}{$4x^2 y^2$}{$2xy^2$}

\vspace{1mm}

\noindent\textbf{\large PHẦN II. TỰ LUẬN (7,5 điểm)}

\noindent\textbf{Câu 1 (2,5 điểm). Thực hiện phép tính:}
\vspace{0.8mm}

\noindent
\begin{tabularx}{\linewidth}{|X|X|}
\hline
\oBaiTap[17]{a)}{$A = (12x^5 y^3 - 8x^4 y^2 + 4x^3 y) : 4x^3 y$} & 
\oBaiTap[17]{b)}{$B = (15x^4 y^3 - 10x^3 y^4) : 5x^2 y^2 - 2xy(x - y)$} \\
\hline
\end{tabularx}


\newpage

% #############################################################################
% ========================= TOÁN 8 - TUẦN 03 - BUỔI 01: MÃ ĐỀ 101 ============
% ========================= MẶT SAU (TRANG 2) =================================
% #############################################################################

\noindent
\begin{tabularx}{\linewidth}{|X|c|}
\hline
\textbf{PHẦN II. TỰ LUẬN (Tiếp theo) - TOÁN 8 TUẦN 03 BUỔI 01} & \textbf{Mã đề: 101} \\
\hline
\end{tabularx}

\vspace{1mm}

\noindent
\begin{tabularx}{\linewidth}{|X|}
\hline
\rule[-2mm]{0pt}{7mm}\textbf{Họ và tên:} \textbf{\large Chu Gia Huy} \hfill \textbf{Mã HS:} \texttt{GIAHUY0503} \quad \\
\hline
\end{tabularx}

\vspace{1.5mm}

\noindent\textbf{\large Câu 2 (2,5 điểm). Rút gọn và tìm $x$:}
\vspace{0.8mm}

\noindent
\begin{tabularx}{\linewidth}{|X|X|}
\hline
\oBaiTap[12]{a)}{Rút gọn và tính giá trị biểu thức:\par
$P = (20x^4 y^3 - 15x^3 y^2) : 5x^2 y^2 - x(4xy - 3)$\par
tại $x = -2$ và $y = 3$.} & 
\oBaiTap[12]{b)}{Tìm số thực $x$, biết:\par
$(12x^3 - 8x^2) : 4x^2 - (x - 1)(2x + 1) + 2x^2 = 7$.} \\
\hline
\end{tabularx}

\vspace{1.8mm}

\noindent\textbf{\large Câu 3 (2,5 điểm). Bài toán thực tế và nâng cao:}
\vspace{0.8mm}

\noindent
\begin{tabularx}{\linewidth}{|X|}
\hline
\vspace{0.5mm}
\textbf{a) (1,5 điểm)} Một hình hộp chữ nhật có thể tích được biểu diễn bởi đa thức $V = 16x^3 y^2 + 24x^2 y^3$ ($\text{cm}^3$). Biết đáy của hình hộp chữ nhật có diện tích là $S = 8x^2 y^2$ ($\text{cm}^2$).\par
1) Hãy viết biểu thức biểu thị chiều cao $h$ của hình hộp chữ nhật theo $x$ và $y$.\par
2) Tính chiều cao của hình hộp chữ nhật khi $x = 4\text{ cm}$ và $y = 3\text{ cm}$.\par
\vspace{0.4mm}
\noindent\textit{Bài giải:}\dongke{9}
\vspace{1mm}
\textbf{b) (1,0 điểm)} Tìm số tự nhiên $n$ để đa thức $A = 12x^{n+2} y^4 - 8x^5 y^{n+1}$ chia hết cho đơn thức $B = 4x^4 y^3$.\par
\vspace{0.4mm}
\noindent\textit{Bài giải:}\dongke{8}
\\ \hline
\end{tabularx}


\newpage

% #############################################################################
% ========================= TOÁN 8 - TUẦN 03 - BUỔI 01: MÃ ĐỀ 102 ============
% ========================= MẶT TRƯỚC (TRANG 3) ===============================
% #############################################################################

\noindent
\begin{tabularx}{\linewidth}{|X|c|}
\hline
\textbf{PHIẾU KIỂM TRA ĐỊNH KỲ TOÁN 8 - TUẦN 03 - BUỔI 01} & \textbf{Mã đề: 102} \\
\textit{Chủ đề: Phép chia đa thức cho đơn thức} & \textit{Thời gian: 45 phút} \\
\hline
\end{tabularx}

\vspace{1mm}

\noindent
\begin{tabularx}{\linewidth}{|X|}
\hline
\rule[-2mm]{0pt}{7.5mm}\textbf{Họ và tên:} \textbf{\large Phan Nhật Vĩnh Phú} \hfill \textbf{Mã HS:} \texttt{VINHPHU2104} \quad \\
\hline
\end{tabularx}

\vspace{1mm}

\noindent
\begin{tabularx}{\linewidth}{|p{3.2cm}|X|}
\hline
\centering\textbf{ĐIỂM SỐ} & \centering\textbf{LỜI NHẬN XÉT CỦA THẦY CÔ} \tabularnewline
\hline
\rule{0pt}{1.3cm} & \tabularnewline
\hline
\end{tabularx}

\vspace{1.2mm}

\noindent\textbf{BẢNG ĐIỀN ĐÁP ÁN TRẮC NGHIỆM (Mỗi câu đúng 0,5 điểm):}
\vspace{0.8mm}

\noindent
\renewcommand{\arraystretch}{1.35}
\begin{tabularx}{\linewidth}{|c|*{5}{>{\centering\arraybackslash}X|}}
\hline
\textbf{Câu} & \textbf{1} & \textbf{2} & \textbf{3} & \textbf{4} & \textbf{5} \\
\hline
\textbf{Đáp án} & & & & & \\
\hline
\end{tabularx}
\renewcommand{\arraystretch}{1.0}

\vspace{1.5mm}

\noindent\textbf{\large PHẦN I. TRẮC NGHIỆM KHÁCH QUAN (2,5 điểm)}
\vspace{1mm}

\cauTN{1}{Kết quả của phép chia đơn thức $18x^6 y^4 : (-3x^3 y^2)$ là:}{$6x^3 y^2$}{$-6x^3 y^2$}{$-6x^2 y^2$}{$-6x^9 y^6$}

\cauTN{2}{Đơn thức $16x^5 y^4$ chia hết cho đơn thức nào dưới đây?}{$5x^6 y^2$}{$4x^4 y^5$}{$-8x^3 y^4$}{$-2x^6 y^6$}

\cauTN{3}{Kết quả phép chia $(15x^5 y^3 - 10x^4 y^2 + 5x^3 y) : 5x^3 y$ là:}{$3x^2 y^2 - 2xy + 1$}{$3x^2 y^2 - 2xy$}{$3x^2 y^2 - 2x + 1$}{$3x^2 y^2 - 2xy + y$}

\cauTN{4}{Đa thức $Q = 8x^4 y^3 - 12x^3 y^2 + 4x^2 y$ chia hết cho đơn thức nào sau đây?}{$4x^3 y$}{$4xy^2$}{$8x^2 y^2$}{$-4x^2 y$}

\cauTN{5}{Biết $(18x^5 y^4 - 12x^4 y^3) : B = 6x^3 y^2 - 4x^2 y$. Đơn thức $B$ là:}{$3xy^2$}{$3x^2 y^2$}{$2x^2 y^2$}{$3x^2 y$}

\vspace{1mm}

\noindent\textbf{\large PHẦN II. TỰ LUẬN (7,5 điểm)}

\noindent\textbf{Câu 1 (2,5 điểm). Thực hiện phép tính:}
\vspace{0.8mm}

\noindent
\begin{tabularx}{\linewidth}{|X|X|}
\hline
\oBaiTap[17]{a)}{$A = (15x^5 y^4 - 9x^4 y^3 + 6x^3 y^2) : 3x^3 y^2$} & 
\oBaiTap[17]{b)}{$B = (18x^4 y^4 - 12x^3 y^3) : 6x^2 y^2 - 3xy(xy - 1)$} \\
\hline
\end{tabularx}


\newpage

% #############################################################################
% ========================= TOÁN 8 - TUẦN 03 - BUỔI 01: MÃ ĐỀ 102 ============
% ========================= MẶT SAU (TRANG 4) =================================
% #############################################################################

\noindent
\begin{tabularx}{\linewidth}{|X|c|}
\hline
\textbf{PHẦN II. TỰ LUẬN (Tiếp theo) - TOÁN 8 TUẦN 03 BUỔI 01} & \textbf{Mã đề: 102} \\
\hline
\end{tabularx}

\vspace{1mm}

\noindent
\begin{tabularx}{\linewidth}{|X|}
\hline
\rule[-2mm]{0pt}{7mm}\textbf{Họ và tên:} \textbf{\large Phan Nhật Vĩnh Phú} \hfill \textbf{Mã HS:} \texttt{VINHPHU2104} \quad \\
\hline
\end{tabularx}

\vspace{1.5mm}

\noindent\textbf{\large Câu 2 (2,5 điểm). Rút gọn và tìm $x$:}
\vspace{0.8mm}

\noindent
\begin{tabularx}{\linewidth}{|X|X|}
\hline
\oBaiTap[12]{a)}{Rút gọn và tính giá trị biểu thức:\par
$Q = (24x^4 y^3 - 18x^3 y^2) : 6x^2 y^2 - x(4xy - 3)$\par
tại $x = 3$ và $y = -2$.} & 
\oBaiTap[12]{b)}{Tìm số thực $x$, biết:\par
$(15x^3 - 10x^2) : 5x^2 - (x - 2)(2x + 1) + 2x^2 = 12$.} \\
\hline
\end{tabularx}

\vspace{1.8mm}

\noindent\textbf{\large Câu 3 (2,5 điểm). Bài toán thực tế và nâng cao:}
\vspace{0.8mm}

\noindent
\begin{tabularx}{\linewidth}{|X|}
\hline
\vspace{0.5mm}
\textbf{a) (1,5 điểm)} Một khối đá dạng hình hộp chữ nhật có thể tích được xác định bởi đa thức $V = 18x^3 y^2 + 27x^2 y^3$ ($\text{cm}^3$). Biết diện tích mặt đáy là $S = 9x^2 y^2$ ($\text{cm}^2$).\par
1) Hãy viết biểu thức biểu thị chiều cao $h$ của khối đá theo $x$ và $y$.\par
2) Tính chiều cao của khối đá khi $x = 3\text{ cm}$ và $y = 2\text{ cm}$.\par
\vspace{0.4mm}
\noindent\textit{Bài giải:}\dongke{9}
\vspace{1mm}
\textbf{b) (1,0 điểm)} Tìm số tự nhiên $n$ để đa thức $A = 15x^{n+3} y^5 - 10x^6 y^{n+2}$ chia hết cho đơn thức $B = 5x^5 y^4$.\par
\vspace{0.4mm}
\noindent\textit{Bài giải:}\dongke{8}
\\ \hline
\end{tabularx}


\newpage

% #############################################################################
% ========================= TOÁN 8 - TUẦN 03 - BUỔI 01: MÃ ĐỀ 103 ============
% ========================= MẶT TRƯỚC (TRANG 5) ===============================
% #############################################################################

\noindent
\begin{tabularx}{\linewidth}{|X|c|}
\hline
\textbf{PHIẾU KIỂM TRA ĐỊNH KỲ TOÁN 8 - TUẦN 03 - BUỔI 01} & \textbf{Mã đề: 103} \\
\textit{Chủ đề: Phép chia đa thức cho đơn thức} & \textit{Thời gian: 45 phút} \\
\hline
\end{tabularx}

\vspace{1mm}

\noindent
\begin{tabularx}{\linewidth}{|X|}
\hline
\rule[-2mm]{0pt}{7.5mm}\textbf{Họ và tên:} \textbf{\large Võ Hồng Phúc} \hfill \textbf{Mã HS:} \texttt{HONGPHUC2201} \quad \\
\hline
\end{tabularx}

\vspace{1mm}

\noindent
\begin{tabularx}{\linewidth}{|p{3.2cm}|X|}
\hline
\centering\textbf{ĐIỂM SỐ} & \centering\textbf{LỜI NHẬN XÉT CỦA THẦY CÔ} \tabularnewline
\hline
\rule{0pt}{1.3cm} & \tabularnewline
\hline
\end{tabularx}

\vspace{1.2mm}

\noindent\textbf{BẢNG ĐIỀN ĐÁP ÁN TRẮC NGHIỆM (Mỗi câu đúng 0,5 điểm):}
\vspace{0.8mm}

\noindent
\renewcommand{\arraystretch}{1.35}
\begin{tabularx}{\linewidth}{|c|*{5}{>{\centering\arraybackslash}X|}}
\hline
\textbf{Câu} & \textbf{1} & \textbf{2} & \textbf{3} & \textbf{4} & \textbf{5} \\
\hline
\textbf{Đáp án} & & & & & \\
\hline
\end{tabularx}
\renewcommand{\arraystretch}{1.0}

\vspace{1.5mm}

\noindent\textbf{\large PHẦN I. TRẮC NGHIỆM KHÁCH QUAN (2,5 điểm)}
\vspace{1mm}

\cauTN{1}{Kết quả của phép chia đơn thức $24x^6 y^4 : 6x^3 y$ là:}{$4x^3 y^3$}{$4x^2 y^3$}{$18x^3 y^3$}{$4x^9 y^5$}

\cauTN{2}{Đơn thức $18x^5 y^3$ chia hết cho đơn thức nào dưới đây?}{$6x^6 y^2$}{$-9x^3 y^3$}{$3x^4 y^4$}{$2x^5 y^5$}

\cauTN{3}{Kết quả phép chia $(18x^4 y^3 - 12x^3 y^2 + 6x^2 y) : 6x^2 y$ là:}{$3x^2 y^2 - 2xy$}{$3x^2 y^2 - 2xy + y$}{$3x^2 y^2 - 2xy + 1$}{$3x^6 y^4 - 2x^5 y^3 + x^4 y^2$}

\cauTN{4}{Đa thức $P = 12x^4 y^2 - 16x^3 y^3 + 8x^2 y$ chia hết cho đơn thức nào sau đây?}{$4x^3 y^2$}{$4x^2 y^2$}{$-4x^2 y$}{$8xy^2$}

\cauTN{5}{Biết $(20x^5 y^4 - 15x^4 y^3) : C = 4x^3 y^2 - 3x^2 y$. Đơn thức $C$ là:}{$5x^2 y$}{$5x^2 y^2$}{$4x^2 y^2$}{$5xy^2$}

\vspace{1mm}

\noindent\textbf{\large PHẦN II. TỰ LUẬN (7,5 điểm)}

\noindent\textbf{Câu 1 (2,5 điểm). Thực hiện phép tính:}
\vspace{0.8mm}

\noindent
\begin{tabularx}{\linewidth}{|X|X|}
\hline
\oBaiTap[17]{a)}{$A = (16x^5 y^4 - 12x^4 y^3 + 8x^3 y^2) : 4x^3 y^2$} & 
\oBaiTap[17]{b)}{$B = (20x^4 y^3 - 15x^3 y^4) : 5x^2 y^2 - 3xy(x - y)$} \\
\hline
\end{tabularx}


\newpage

% #############################################################################
% ========================= TOÁN 8 - TUẦN 03 - BUỔI 01: MÃ ĐỀ 103 ============
% ========================= MẶT SAU (TRANG 6) =================================
% #############################################################################

\noindent
\begin{tabularx}{\linewidth}{|X|c|}
\hline
\textbf{PHẦN II. TỰ LUẬN (Tiếp theo) - TOÁN 8 TUẦN 03 BUỔI 01} & \textbf{Mã đề: 103} \\
\hline
\end{tabularx}

\vspace{1mm}

\noindent
\begin{tabularx}{\linewidth}{|X|}
\hline
\rule[-2mm]{0pt}{7mm}\textbf{Họ và tên:} \textbf{\large Võ Hồng Phúc} \hfill \textbf{Mã HS:} \texttt{HONGPHUC2201} \quad \\
\hline
\end{tabularx}

\vspace{1.5mm}

\noindent\textbf{\large Câu 2 (2,5 điểm). Rút gọn và tìm $x$:}
\vspace{0.8mm}

\noindent
\begin{tabularx}{\linewidth}{|X|X|}
\hline
\oBaiTap[12]{a)}{Rút gọn và tính giá trị biểu thức:\par
$Q = (28x^4 y^3 - 21x^3 y^2) : 7x^2 y^2 - x(4xy - 3)$\par
tại $x = 2$ và $y = -3$.} & 
\oBaiTap[12]{b)}{Tìm số thực $x$, biết:\par
$(18x^3 - 12x^2) : 6x^2 - (x - 1)(3x + 1) + 3x^2 = 8$.} \\
\hline
\end{tabularx}

\vspace{1.8mm}

\noindent\textbf{\large Câu 3 (2,5 điểm). Bài toán thực tế và nâng cao:}
\vspace{0.8mm}

\noindent
\begin{tabularx}{\linewidth}{|X|}
\hline
\vspace{0.5mm}
\textbf{a) (1,5 điểm)} Một bể nước dạng hình hộp chữ nhật có thể tích được biểu diễn bởi đa thức $V = 20x^3 y^2 + 30x^2 y^3$ ($\text{cm}^3$). Biết đáy bể có diện tích là $S = 10x^2 y^2$ ($\text{cm}^2$).\par
1) Hãy viết biểu thức biểu thị chiều cao $h$ của bể nước theo $x$ và $y$.\par
2) Tính chiều cao của bể nước khi $x = 5\text{ cm}$ và $y = 2\text{ cm}$.\par
\vspace{0.4mm}
\noindent\textit{Bài giải:}\dongke{9}
\vspace{1mm}
\textbf{b) (1,0 điểm)} Tìm số tự nhiên $n$ để đa thức $A = 16x^{n+2} y^5 - 12x^5 y^{n+1}$ chia hết cho đơn thức $B = 4x^4 y^3$.\par
\vspace{0.4mm}
\noindent\textit{Bài giải:}\dongke{8}
\\ \hline
\end{tabularx}


\newpage

% #############################################################################
% ========================= TOÁN 8 - TUẦN 03 - BUỔI 01: MÃ ĐỀ 104 ============
% ========================= MẶT TRƯỚC (TRANG 7) ===============================
% #############################################################################

\noindent
\begin{tabularx}{\linewidth}{|X|c|}
\hline
\textbf{PHIẾU KIỂM TRA ĐỊNH KỲ TOÁN 8 - TUẦN 03 - BUỔI 01} & \textbf{Mã đề: 104} \\
\textit{Chủ đề: Phép chia đa thức cho đơn thức} & \textit{Thời gian: 45 phút} \\
\hline
\end{tabularx}

\vspace{1mm}

\noindent
\begin{tabularx}{\linewidth}{|X|}
\hline
\rule[-2mm]{0pt}{7.5mm}\textbf{Họ và tên:} \textbf{\large Nguyễn Văn Thông} \hfill \textbf{Mã HS:} \texttt{VANTHONG1902} \quad \\
\hline
\end{tabularx}

\vspace{1mm}

\noindent
\begin{tabularx}{\linewidth}{|p{3.2cm}|X|}
\hline
\centering\textbf{ĐIỂM SỐ} & \centering\textbf{LỜI NHẬN XÉT CỦA THẦY CÔ} \tabularnewline
\hline
\rule{0pt}{1.3cm} & \tabularnewline
\hline
\end{tabularx}

\vspace{1.2mm}

\noindent\textbf{BẢNG ĐIỀN ĐÁP ÁN TRẮC NGHIỆM (Mỗi câu đúng 0,5 điểm):}
\vspace{0.8mm}

\noindent
\renewcommand{\arraystretch}{1.35}
\begin{tabularx}{\linewidth}{|c|*{5}{>{\centering\arraybackslash}X|}}
\hline
\textbf{Câu} & \textbf{1} & \textbf{2} & \textbf{3} & \textbf{4} & \textbf{5} \\
\hline
\textbf{Đáp án} & & & & & \\
\hline
\end{tabularx}
\renewcommand{\arraystretch}{1.0}

\vspace{1.5mm}

\noindent\textbf{\large PHẦN I. TRẮC NGHIỆM KHÁCH QUAN (2,5 điểm)}
\vspace{1mm}

\cauTN{1}{Kết quả của phép chia đơn thức $21x^6 y^5 : (-7x^2 y^3)$ là:}{$-3x^4 y^2$}{$3x^4 y^2$}{$-3x^3 y^2$}{$-3x^8 y^8$}

\cauTN{2}{Đơn thức $20x^5 y^4$ chia hết cho đơn thức nào dưới đây?}{$4x^6 y^3$}{$5x^4 y^5$}{$-4x^4 y^3$}{$-10x^5 y^5$}

\cauTN{3}{Kết quả phép chia $(21x^5 y^4 - 14x^4 y^3 + 7x^3 y^2) : 7x^3 y^2$ là:}{$3x^2 y^2 - 2xy$}{$3x^2 y^2 - 2xy + 1$}{$3x^2 y^2 - 2x + 1$}{$3x^2 y^2 - 2xy + y$}

\cauTN{4}{Đa thức $Q = 15x^4 y^3 - 20x^3 y^2 + 10x^2 y$ chia hết cho đơn thức nào sau đây?}{$-5x^2 y$}{$5x^3 y$}{$5xy^2$}{$10x^2 y^2$}

\cauTN{5}{Biết $(24x^5 y^4 - 18x^4 y^3) : D = 4x^3 y^2 - 3x^2 y$. Đơn thức $D$ là:}{$6xy^2$}{$6x^2 y$}{$6x^2 y^2$}{$4x^2 y^2$}

\vspace{1mm}

\noindent\textbf{\large PHẦN II. TỰ LUẬN (7,5 điểm)}

\noindent\textbf{Câu 1 (2,5 điểm). Thực hiện phép tính:}
\vspace{0.8mm}

\noindent
\begin{tabularx}{\linewidth}{|X|X|}
\hline
\oBaiTap[17]{a)}{$A = (18x^5 y^4 - 12x^4 y^3 + 6x^3 y^2) : 6x^3 y^2$} & 
\oBaiTap[17]{b)}{$B = (24x^4 y^4 - 16x^3 y^3) : 8x^2 y^2 - 2xy(xy - 1)$} \\
\hline
\end{tabularx}


\newpage

% #############################################################################
% ========================= TOÁN 8 - TUẦN 03 - BUỔI 01: MÃ ĐỀ 104 ============
% ========================= MẶT SAU (TRANG 8) =================================
% #############################################################################

\noindent
\begin{tabularx}{\linewidth}{|X|c|}
\hline
\textbf{PHẦN II. TỰ LUẬN (Tiếp theo) - TOÁN 8 TUẦN 03 BUỔI 01} & \textbf{Mã đề: 104} \\
\hline
\end{tabularx}

\vspace{1mm}

\noindent
\begin{tabularx}{\linewidth}{|X|}
\hline
\rule[-2mm]{0pt}{7mm}\textbf{Họ và tên:} \textbf{\large Nguyễn Văn Thông} \hfill \textbf{Mã HS:} \texttt{VANTHONG1902} \quad \\
\hline
\end{tabularx}

\vspace{1.5mm}

\noindent\textbf{\large Câu 2 (2,5 điểm). Rút gọn và tìm $x$:}
\vspace{0.8mm}

\noindent
\begin{tabularx}{\linewidth}{|X|X|}
\hline
\oBaiTap[12]{a)}{Rút gọn và tính giá trị biểu thức:\par
$Q = (32x^4 y^3 - 24x^3 y^2) : 8x^2 y^2 - x(4xy - 3)$\par
tại $x = -2$ và $y = 2$.} & 
\oBaiTap[12]{b)}{Tìm số thực $x$, biết:\par
$(20x^3 - 15x^2) : 5x^2 - (x - 2)(3x + 1) + 3x^2 = 14$.} \\
\hline
\end{tabularx}

\vspace{1.8mm}

\noindent\textbf{\large Câu 3 (2,5 điểm). Bài toán thực tế và nâng cao:}
\vspace{0.8mm}

\noindent
\begin{tabularx}{\linewidth}{|X|}
\hline
\vspace{0.5mm}
\textbf{a) (1,5 điểm)} Một thùng chứa hình hộp chữ nhật có thể tích là $V = 24x^3 y^2 + 36x^2 y^3$ ($\text{cm}^3$). Biết diện tích đáy là $S = 12x^2 y^2$ ($\text{cm}^2$).\par
1) Hãy viết biểu thức biểu thị chiều cao $h$ của thùng chứa theo $x$ và $y$.\par
2) Tính chiều cao của thùng chứa khi $x = 2\text{ cm}$ và $y = 3\text{ cm}$.\par
\vspace{0.4mm}
\noindent\textit{Bài giải:}\dongke{9}
\vspace{1mm}
\textbf{b) (1,0 điểm)} Tìm số tự nhiên $n$ để đa thức $A = 20x^{n+3} y^5 - 15x^6 y^{n+2}$ chia hết cho đơn thức $B = 5x^5 y^4$.\par
\vspace{0.4mm}
\noindent\textit{Bài giải:}\dongke{8}
\\ \hline
\end{tabularx}



\end{document}
